\title{Parameter-Efficient Orthogonal Finetuning via Factorization}

\begin{abstract}
With the growing prevalence of large foundation models, the costs associated with training these models from the ground up have become increasingly prohibitive. Consequently, the ability to effectively adapt these sophisticated models to various downstream applications has gained importance. In this study, we investigate a systematic finetuning approach known as Orthogonal Finetuning (OFT) for adapting models to new tasks. While OFT exhibits strong generalization, its reliance on high-dimensional orthogonal matrices results in a substantial count of trainable parameters. To mitigate this, we analyze OFT through the lens of information transmission, extracting several fundamental criteria essential for enhancing parameter efficiency. Drawing inspiration from the Cooley-Tukey fast Fourier transform algorithm, which achieves efficient information transmission, we devise a more parameter-efficient orthogonal representation leveraging butterfly structures. Incorporating this parameterization into OFT yields a new, parameter-efficient finetuning method referred to as Orthogonal Butterfly~(BOFT). BOFT encompasses OFT as a particular instance, thus presenting a broader orthogonal finetuning paradigm. Finally, we conduct comprehensive empirical experiments on transferring large vision transformers, language models, and text-to-image diffusion models across diverse vision and language downstream tasks.
\end{abstract}

\section{Introduction}

Recent advances such as ChatGPT~\cite{brown2020language,chen2021evaluating} and Stable Diffusion~\cite{rombach2022high} illustrate the impressive generalization capabilities that large foundation models possess. However, as model performance rapidly advances, the parameter scales have soared

(\emph{e.g.}, GPT-3 contains approximately 175B parameters), posing considerable obstacles to training such foundation models from scratch. Accordingly, significant advancements now depend on strategies for repurposing existing foundation models for novel tasks, without initiating training anew. In other words, developing parameter- and computation-efficient means to adapt foundation models is crucial. Predominantly, three strategies exist for efficient task adaptation: \emph{model finetuning}~\cite{hulora2022,zaken2022bitfit,guo2020parameter,raffel2020exploring,qiu2023controlling,zhang2022adaptive,chavan2023one}, which finetunes a portion of parameters while retaining the architecture; \emph{adapter tuning}~\cite{rebuffi2017learning,houlsby2019parameter,pfeiffer2020adapterfusion,lin2020exploring,he2021towards}, which augments the model with additional trainable components; and \emph{prompt tuning}~\cite{li2021prefix,lester2021power}, which learns supplementary prefix tokens prepended to the input. Of these, model finetuning stands out due to its conceptual simplicity, robust performance, and, notably, the absence of increased inference latency.

The core concept underpinning model finetuning is to maintain a degree of similarity between the pretrained model and its finetuned counterpart, according to certain metrics, thereby retaining the knowledge acquired during pretraining. Standard finetuning techniques typically employ a reduced learning rate—a heuristic choice that restricts divergence between the pretrained and finetuned models. Given the inherent structure of weight matrices, more systematic strategies aim to conserve relational attributes among these matrices, such as the pairwise angular relationships between neurons. This rationale motivates the development of Orthogonal Finetuning~(OFT)~\cite{qiu2023controlling}, a novel finetuning paradigm where neurons within a layer are jointly acted upon by a shared orthogonal matrix, thus rigorously preserving pairwise neuron angles across the finetuning trajectory. OFT has yielded strong results in terms of generalization and convergence when applied to finetuning text-to-image diffusion models~\cite{qiu2023controlling}. Nonetheless, the necessity of learning full orthogonal matrices results in a large parameter count due to their high dimensionality. To mitigate this, OFT imposes a block-diagonal configuration, reducing the parameter overhead. However, this parameter-saving scheme brings certain drawbacks—namely, the imposed fixed sparsity pattern and the independent operation of orthogonal transformations within each block. Such a structure, while efficient, may introduce inductive biases (\emph{e.g.}, the block-diagonal orthogonal matrices reduce model expressivity and are unable to represent conventional linear transformations), and crucially, optimal determination of the sparsity structure remains unresolved.

A central challenge in this context lies in constructing a dense orthogonal matrix without incurring excessive parameter overhead. Traditionally, a dense orthogonal matrix of dimensionality $d$ requires $\mathcal{O}(d^2)$ parameters, which is computationally demanding. In contrast, we introduce an innovative strategy that constructs such a matrix by composing several sparse, factorized matrices. This method is motivated by the principle of reducing the number of trainable parameters by trading increased computation during training for decreased storage. Since the orthogonal matrix is expressed as a product of sparse matrices, these matrix multiplications must be repeated throughout each iteration of training. To conceptualize this factorization process, we frame the generation of a dense orthogonal matrix as a problem of information flow. In this formulation, creating a dense orthogonal matrix by sequentially multiplying sparse matrices is analogous to transmitting information across a grid-like graph structure. This lens allows for a rich design space in sparse matrix factorizations, providing flexibility to restrict trainable parameters yet maintain the ability to represent dense matrices effectively.

To further enhance parameter efficiency within the information transmission paradigm, we draw inspiration from the butterfly architecture found in the Cooley-Tukey fast Fourier transform algorithm~\cite{cooley1965algorithm}, known for its highly efficient information exchange between nodes~\cite{klasing1994broadcasting}. The butterfly network naturally leads to a method of sparse matrix factorization, referred to as butterfly factorization. Utilizing this structure, we can generate a dense $d$-dimensional matrix as the product of $\mathcal{O}(\log d)$ sparse matrices, each containing only $\mathcal{O}(d)$ nonzero elements. By ensuring each sparse matrix in the product is orthogonal, we obtain a parameterization requiring just $\mathcal{O}(d\log d)$ parameters—a drastic reduction compared to the standard dense formulation. Building on this efficient orthogonal construction, we introduce Orthogonal Butterfly~(BOFT), a new approach for parameter-efficient finetuning. BOFT generalizes OFT, including it as a particular instance, thereby offering a comprehensive framework for orthogonal finetuning. Notably, both the block-diagonal approach and the butterfly approach share a core property—sparsity—which allows orthogonal matrices to maintain expressiveness while minimizing the number of parameters that need to be learned. Our methodology can be fruitfully compared to the low-rank decompositions utilized in LoRA~\cite{hulora2022}, as both low-rank and sparse matrices serve as prominent classes of structured matrices~\cite{candes2011robust} designed for efficient parameter utilization.

In contrast to the block-diagonal structure adopted by OFT, which mediates between expressivity and regularity, BOFT leverages the butterfly structure to enable a more continuous transition between matrices from the full orthogonal group (for expressivity) and identity matrices (for regularity). This facilitates the identification of a more compact hypothesis space within the orthogonal group for downstream applications. Given the extensive adoption of the butterfly structure in numerous fast linear operations—such as the discrete Fourier and discrete cosine transforms—we posit that our structurally constrained approach to parameter efficiency naturally introduces an implicit inductive bias, which enhances generalization and curbs overfitting. This reasoning is grounded in the fact that the butterfly structure inherently accommodates a wide range of classical linear transforms, a property absent in the block-diagonal construction used by OFT. The main contributions of our work are summarized below:

\begin{itemize}[leftmargin=*,nosep]
    \item We introduce a novel information transmission framework for studying parameter efficiency in orthogonal finetuning. This framework recasts the challenge of designing parameter-efficient dense orthogonal matrices as an information flow problem on a grid-based graph structure.
    \item Building on the butterfly patterns from the Cooley-Tukey algorithm, we introduce Orthogonal Butterfly, an orthogonal finetuning method with high parameter efficiency, formulated within our information transmission paradigm.
    \item We offer several theoretical explanations for the substantial reduction in trainable parameters achieved by BOFT, along with maintained expressivity and stable optimization during training. Owing to the matrix factorization technique, BOFT also naturally enables a form of weight interpolation (refer to Figure~\ref{fig:boft_interpolation}).
    \item We pioneer the application of orthogonal finetuning in a broad spectrum of tasks extending beyond controllable text-to-image generation~\cite{qiu2023controlling}, highlighting its promise as a versatile model finetuning strategy. To demonstrate this, we utilize BOFT across diverse adaptation scenarios, from computer vision to natural language processing, where it surpasses leading state-of-the-art baselines, thereby corroborating its advanced parameter efficiency and generalization performance.
\end{itemize}

\textbf{Parameter-efficient finetuning (PEFT)}. With the rapid expansion and increasing capability of foundation models (\emph{e.g.}, \cite{brown2020language,radford2021learning,rombach2022high,kirillov2023segment}), there is heightened interest in developing approaches to finetune these models for specific downstream tasks while optimizing parameter efficiency~\cite{treviso2023efficient,ding2023parameter,lialin2023scaling}. Numerous PEFT strategies have been proposed~\cite{houlsby2019parameter,aghajanyan2020intrinsic,hulora2022,edalati2022krona,wang2022adamix,gheini2021cross,zaken2022bitfit,guo2020parameter,sung2021training,ansell2022composable,lester2021power,li2021prefix,vu2022spot,he2021towards,mao2021unipelt,karimi2021compacter,liu2022few,sung2022lst,chen2022parameter,jia2022visual,chen2022adaptformer,zhang2022neural,jie2023fact,lian2022scaling,luo2023towards}; of these, methods leveraging reparameterization~\cite{aghajanyan2020intrinsic,hulora2022,edalati2022krona} are particularly germane to our study. The LoRA framework~\cite{hulora2022} introduces an additional term—a product of two low-rank matrices—on top of the pretrained weight matrix, enabling notable improvements on various natural language benchmarks. While LoRA assigns a fixed rank to all inserted matrices, several enhancements \cite{zhang2022adaptive,valipour2022dylora,zhang2023increlora} have been put forward to adaptively vary the rank across different layers, yielding a more judicious distribution of the parameter budget. The straightforwardness of low-rank reparameterization has propelled its widespread adoption \cite{dettmers2023qlora,zi2023delta,chavan2023one}. Drawing from the use of hyperspherical energy to explain generalization properties~\cite{liu2018learning,liu2021orthogonal}, \cite{qiu2023controlling} introduces orthogonal finetuning (OFT)—an alternative technique for adapting text-to-image diffusion models. OFT works by training an orthogonal transformation on the activations within each layer, and consistently surpasses LoRA in terms of generalization performance and training stability. Nonetheless, this approach tends to increase the count of trainable parameters relative to LoRA, making further reductions in OFT’s parameter footprint an appealing objective. Additionally, the utility of OFT in adaptation scenarios outside text-to-image diffusion remains unverified. BOFT addresses the parameter inefficiency of OFT by applying butterfly factorization, thereby facilitating the application of orthogonal finetuning to a broad range of adaptation challenges.

\textbf{Butterfly structures}. The radix-2 Cooley-Tukey algorithm efficiently computes the $N$-point discrete Fourier transform by recursively partitioning it into two $\frac{N}{2}$-point discrete Fourier transforms, resulting in a characteristic butterfly computation pattern that can be expressed as a sequence of sparse matrix products (often termed butterfly matrices). Butterfly matrices have been employed for parameterizing orthogonal matrices in order to bypass pivoting steps in Gaussian elimination and to enhance computational efficiency~\cite{parker1995random}; they also support stable training of recurrent neural networks~\cite{jing2017tunable} and are valuable in kernel function approximation~\cite{munkhoeva2018quadrature}. Works such as~\cite{dao2019learning,dao2019kaleidoscope} focus on learning efficient linear transforms using butterfly-based parameterizations, while \cite{chen2022pixelated,dao2022monarch} exploit butterfly matrices to streamline neural network training. Furthermore, butterfly structures have been leveraged in accelerating matrix-vector operations~\cite{michielssen1996multilevel,o2010algorithm}, constructing data-sparse matrix approximations~\cite{li2015butterfly}, and optimizing network communications~\cite{klasing1994broadcasting,li2003linear,rajasingh2016transmission}. In our work, distinct from these applications, butterfly structures are explored as a mechanism to improve the parameter efficiency of OFT.

\section{An Information Transmission View on Orthogonal Finetuning}

We begin by outlining the foundational concepts of OFT. In order to adapt the pretrained weight matrix for finetuning, OFT reformulates the updated weight matrix as the product of a learnable orthogonal matrix and the static pretrained weight matrix. Unlike LoRA, which modifies weights via the addition of a low-rank matrix, OFT employs a multiplicative approach using an orthogonal matrix. LoRA’s parameter-efficiency stems from its use of a low-rank design, while the original OFT achieves efficiency through a block-diagonal orthogonal construction~\cite{qiu2023controlling}; decreasing the size of each diagonal block improves OFT’s parameter efficiency. Figure~\ref{fig:comp_lora} illustrates a direct comparison. The rationale for utilizing orthogonal transformations in finetuning is to maintain the pairwise angles between neuron vectors~\cite{liu2018learning,liu2021orthogonal,qiu2023controlling}, thereby preserving the semantic information acquired during pretraining. 

More specifically, OFT learns an orthogonal matrix $\bm{R}\in\mathbb{R}^{d\times d}$ in conjunction with a pretrained linear layer $\bm{W}^0\in\mathbb{R}^{d\times n}$, changing the forward computation from $\bm{z}=(\bm{W}^0)^\top\bm{x}$ to $\bm{z}=(\bm{R}\bm{W}^0)^\top\bm{x}$, where $\bm{x}\in\mathbb{R}^{d}$ represents the input vector and $\bm{z}\in\mathbb{R}^{n}$ the output. For zero initialization, OFT sets $\bm{R}$ to the identity matrix at the start. In order to preserve the orthogonality of $\bm{R}$ throughout the finetuning procedure, we adopt the Cayley parameterization as in \cite{qiu2023controlling,liu2021orthogonal}, specifically $\bm{R}=(\bm{I}+\bm{Q})(\bm{I}-\bm{Q})^{-1}$, where $\bm{Q}$ is a skew-symmetric matrix satisfying $\bm{Q}=-\bm{Q}^\top$. To enhance parameter-efficiency, a block-diagonal structure is implemented by representing $\bm{R}$ as $\text{diag}(\bm{R}_1,\bm{R}_2,\cdots,\bm{R}_r)$, with each $\bm{R}_i\in\mathcal{R}^{b\times b}$ being a compact orthogonal block, and $br=d$. 

However, this block-diagonal design, while parameter-efficient, imposes a structural assumption: each neuron's dimensionality (the components of a column in $\bm{W}^0$) is partitioned into $r$ subsets, each of which is independently transformed via distinct orthogonal matrices. Although this configuration has shown practical success, the underlying logic of dividing a neuron's dimensions based purely on index order is unconvincing, highlighting the potential advantages of using dense orthogonal matrices. This leads to an important question: \emph{Is it possible to devise a dense orthogonal matrix that does not sacrifice parameter-efficiency?}

To tackle this issue, we suggest constructing a dense orthogonal matrix by multiplying several sparse orthogonal matrices together. To offer a comprehensive yet accessible approach to analyze the sparsity structures in orthogonal matrix decomposition, we conceptualize the generation of a dense orthogonal matrix within OFT as an instance of information transmission. More precisely, forming a dense matrix $\bm{R}\in\mathbb{R}^{d\times d}$ through the product of $m$ square matrices, expressed as $\thickmuskip=2mu \medmuskip=2mu \bm{R}=\bm{B}_m\bm{B}_{m-1}\cdots\bm{B}_1$, can be interpreted as relaying information across a $d\times (m+1)$ grid, as depicted in Figure~\ref{fig:info_trans}. The rationale for adopting the information transmission analogy arises from the idea that a dense $d$-by-$d$ matrix naturally represents complete connections between two sets of $d$ nodes. In this context, for matrix $\bm{R}$, a zero entry $R_{ij}$ indicates the absence of a path for information flow from node $j$ to node $i$, whereas a nonzero $R_{ij}$ means such transmission is feasible. Thus, describing the dense matrix $\bm{R}$ as a sequence of products $\bm{B}_m\bm{B}_{m-1}\cdots\bm{B}_1$ can also be seen as carrying out a series of information exchanges dictated by the graph structures corresponding to each $\bm{B}_i$. The flow of information starts with the transformation imposed by $\bm{B}_1$ and concludes with $\bm{B}_n$. As a specific case illustrated in Figure~\ref{fig:info_trans}, we look at the decomposition $\bm{R}=\bm{B}_5\bm{B}_4\bm{B}_3\bm{B}_2\bm{B}_1$ with the associated sparsity patterns and resultant graph shown. The graph in Figure~\ref{fig:info_trans} reflects an expanded view of the matrix multiplication. Here, each $\bm{B}_i$ serves as a connectivity map from nodes in level $i$ to those in level $i+1$. Explicitly, the $(j_1,j_2)$ entry of $\bm{B}_i$ specifies whether a directed edge exists from node $j_2$ on level $i$ to node $j_1$ on level $i+1$ (a zero means no connection). For the product $\bm{B}_5\bm{B}_4\bm{B}_3\bm{B}_2\bm{B}_1$ to form a dense matrix, information must be able to traverse from any node at the first level to every node at the sixth level. Examining $\bm{R}=\bm{B}_4\bm{B}_3\bm{B}_2\bm{B}_1$—linking nodes from level one to level five—shows that information from node $1$ cannot reach node $3$, which means that the corresponding sparsity pattern of $\bm{B}_4\bm{B}_3\bm{B}_2\bm{B}_1$ includes a zero at position $(3,1)$. This correspondence between the sparsity structure and the induced graph holds for shorter products as well, such as $\bm{R}=\bm{B}_3\bm{B}_2\bm{B}_1$ or $\bm{R}=\bm{B}_2\bm{B}_1$. In general terms, for $\bm{R}\in\mathbb{R}^{d\times d}$ to be dense, the collection of $m$ matrices $\bm{B}_m,\cdots,\bm{B}_1$ must together describe directed edges on a $d\times (m+1)$ grid—where each directed edge links nodes in neighboring levels (i.e., columns)—in such a way that each node at the initial level is able to transmit information to every node at the terminal level $(m+1)$.

Adopting the information transmission perspective provides a clearer rationale for interpreting the sparsity structure inherent in the factorization matrices used in OFT. As illustrated in Figure~\ref{fig:blk_diag}, the original OFT imposes a block-diagonal pattern on $\bm{R}$. While this structure successfully decreases the number of learnable parameters, it inherently precludes the realization of a densely populated matrix $\bm{R}$. Our primary objective is thus to construct a dense orthogonal matrix by composing $m$ sparse orthogonal matrices, striving to utilize the smallest set of effective trainable parameters. Framed within the information transmission context, two core requirements emerge: \emph{(i)} \textbf{dense connectivity}, ensuring every node at the initial level can reach all nodes at the final level via at least one path, and \emph{(ii)} \textbf{minimal free edges}, meaning the network of edges must be kept as compact as possible given the orthogonality constraint. Orthogonality enforces nuanced restrictions on the allowable edges between subsequent levels. Specifically, for each $\bm{B}_i$ to achieve full-rank—essential for orthogonality—it is necessary to establish $d$ edges forming a bijective mapping between nodes at level $i$ and nodes at level $(i=1)$, thereby requiring no fewer than $d$ connections between adjacent layers (for instance, as demonstrated with four edges in Figure~\ref{fig:info_trans}). These $d$ connections are indispensable for orthogonality and do not constitute learnable parameters, as in a $d\times d$ orthogonal matrix with $d$ non-zero elements, such values are constrained to be $\pm 1$. Since orthogonal matrices possess a lower parameter count than fully dense matrices, enforcing orthogonality induces complex dependencies among edges. For instance, in the case of a $2\times 2$ orthogonal matrix, setting the $(1,1)$ entry to zero also necessitates the $(2,2)$ entry being zero—implying that the elimination of one edge may necessitate the loss of another. Consequently, for any admissible edge configuration, the orthogonality condition might either introduce or eliminate certain edges. Examining the support pattern—the locations of non-zero elements—of the resultant orthogonal matrix, the information transmission framework proves especially apt for analyzing OFT, given that the focus is solely on which entries of $\bm{R}$ are non-zero and not on their explicit values.

Establishing a straightforward fully connected network between two layers involves $\mathcal{O}(d^2)$ connections, corresponding to a dense orthogonal matrix, which translates to $d^2-d$ distinct edges (for $d=4$, this equals 12 edges). As illustrated in Figure~\ref{fig:info_trans}, an example of an achievable matrix factorization requires only $10$ edges, notably fewer than those needed for a dense orthogonal matrix. This representation enables examination of OFT's parameter efficiency through a graph-theoretic lens, facilitating the design of practical factorizations within this structure. Our design draws from the architecture known as butterfly graphs~\cite{cooley1965algorithm}—a structure originating in the Cooley-Tukey algorithm—which effectively links $d$ input and $d$ output nodes using merely $\mathcal{O}(d\log d)$ edges. To illustrate, the topology in Figure~\ref{fig:info_trans} achieves full connectivity using $10$ edges, whereas a butterfly network attains this using just $8$ edges. In the following discussion, we present the ways in which leveraging the butterfly structure boosts parameter efficiency.

\section{Orthogonal Parameterization by Butterfly Factorization}

The butterfly network, first utilized within the Cooley-Tukey algorithm for efficient computation of the fast Fourier transform, allows local adjustments in the frequency domain to result in global effects in the spatial domain. This idea parallels our challenge in information propagation: ensuring each node in an initial layer communicates information to every node in the terminal layer. Beyond signal processing, the butterfly structure has also become a favored topology in computer networks~\cite{solihin2015fundamentals,li2011linear} to promote efficient data exchange. Assuming $k\geq 2$ where $k$ is a power of $2$, we formalize the \emph{butterfly factor} $\bm{B}^F(k)$ as:

\begin{equation}\label{eq:but_factor}
\begin{aligned}
    \bm{B}^F(k)=\begin{bmatrix}
    \text{diag}(\bm{d}_1) & \text{diag}(\bm{d}_2) \\
    \text{diag}(\bm{d}_3) & \text{diag}(\bm{d}_4)
    \end{bmatrix}\in\mathbb{R}^{k\times k},
\end{aligned}
\end{equation}

where each $\bm{d}_i\in\mathbb{R}^{\frac{k}{2}}$ for $i=1,2,3,4$ are vectors. For $d=2^N$, we introduce the recursive construction of the $d$-dimensional \emph{butterfly matrix} $\bm{B}(d)\in\mathbb{R}^{d\times d}$:

\begin{equation}
\begin{aligned}
    \!\!\bm{B}(d)=\tilde{\bm{B}}(d,d)\!\cdot\!\begin{bmatrix}
    \bm{B}_1(\frac{d}{2})\!\!\!\!\! & \bm{0} \\
    \bm{0}\!\!\!\!\! &\bm{B}_2(\frac{d}{2})
    \end{bmatrix}= \tilde{\bm{B}}(d,d)\tilde{\bm{B}}(d,\frac{d}{2})\cdots\tilde{\bm{B}}(d,2),
\end{aligned}
\end{equation}

Here, $\bm{B}_1(\frac{d}{2})$ and $\bm{B}_2(\frac{d}{2})$ denote two butterfly matrices, each of dimension $\frac{d}{2}$. We introduce the notion of a \emph{butterfly component} as $\thickmuskip=2mu \medmuskip=2mu \tilde{\bm{B}}(d,k)=\text{diag}(\bm{B}^F_1(k),\cdots,\bm{B}^F_{d/k}(k))$, representing a block-diagonal matrix of dimensions $d\times d$ with individual block size $k$. Here, each $\bm{B}^F_i(k)$ for all $i$ refers to a butterfly factor as specified in Equation~\ref{eq:but_factor}. This formulation enables the use of the butterfly matrix framework to represent an orthogonal matrix. The main requirement is that each multiplicative factor $\tilde{\bm{B}}(d,k)$ in the full butterfly matrix $\bm{B}(d)$ must itself be orthogonal. Considering first the case where $k=2$, the block-diagonal structure $\tilde{\bm{B}}(d,2)$ can be straightforwardly constrained to be orthogonal by parameterizing each $2\times 2$ block using a Cayley transform (equivalent to a $2$-dimensional rotation), following the methodology in \cite{liu2021orthogonal,qiu2023controlling}. The sparsity structure (the arrangement of nonzero entries) in any butterfly component is, in essence, a permutation of the sparsity pattern found in $\tilde{\bm{B}}(d,2)$, enabling all butterfly components to be efficiently represented as orthogonal matrices. Thus, this scheme realizes an efficient orthogonal matrix parameterization built from numerous $2\times 2$ orthogonal blocks. In line with the approach of \cite{chen2022pixelated}, we generalize the butterfly matrix to introduce the block butterfly component $\tilde{\bm{B}}^b(d,k)$, where each scalar entry in $\bm{d}_i$ is replaced by a $b\times b$ matrix. To ensure that the block butterfly component $\tilde{\bm{B}}^b(d,2)$ remains orthogonal, each $2b\times 2b$ block is constructed as an orthogonal matrix. Similarly, for the remaining components $\tilde{\bm{B}}^b(d,k)$ with $k>2$, their nonzero structure is given by a blockwise permutation of the sparsity of $\tilde{\bm{B}}^b(d,2)$, thus they can be parameterized in an analogous orthogonal manner. Collectively, these developments yield the forward computation in BOFT as follows:

\begin{equation}
    \begin{aligned}\nonumber
    \bm{z} = \big(\bm{R}(m,b)\cdot\bm{W}^0\big)^\top\bm{x},~~~~\text{s.t.}~\bigg\{\bm{R}(m,b)=\prod_{i=1}^m \tilde{\bm{B}}^b_{(d,i)}~~\&~~\underbrace{\big(\tilde{\bm{B}}^b_{(d,j)}\big)^\top\tilde{\bm{B}}^b_{(d,j)} = \tilde{\bm{B}}^b_{(d,j)}\big(\tilde{\bm{B}}^b_{(d,j)}\big)^\top = \bm{I}_d}_{\forall j\in [1, m]}\bigg\},
    \end{aligned}
\end{equation}

where, for notational brevity, we define $\tilde{\bm{B}}^b(d,2^{m - i+1})$ as $\tilde{\bm{B}}^b_{(d,i)}$, and $\bm{I}_d$ designates the $d$-dimensional identity matrix. The orthogonal matrix $\bm{R}(m,b) \in \mathbb{R}^{d \times d}$ is assembled by multiplying several orthogonal butterfly matrices. For ease of notation, we refer to BOFT employing $\bm{R}(m,\frac{b}{2})$ as BOFT($m$,$b$), with $b\geq 2$. Specifically, when $m=1$, BOFT($1$,$b$) coincides with the block-diagonal OFT~\cite{qiu2023controlling} featuring block size $b$. If $b=d$, then BOFT($1$,$d$) recovers the standard OFT~\cite{qiu2023controlling} using a full, unconstrained orthogonal matrix. Using BOFT($\log \frac{2d}{b}$,$b$), the method constructs $\bm{R}$ as a block butterfly matrix $\bm{B}^{\frac{b}{2}}(d)$, yielding a fully dense orthogonal transformation $\bm{R}$. Broadly, BOFT($m$,$b$) requires $\frac{1}{2}(b-1)dm$ learnable parameters to adapt a $d\times n$ linear transformation. Setting $m=\log d$ and $b=2$ (the butterfly case) reduces parameter complexity to $\mathcal{O}(d\log d)$. By contrast, the traditional OFT based on a full orthogonal matrix incurs $\mathcal{O}(d^2)$ parameters, and block-diagonal OFT with $r$ blocks uses $\mathcal{O}(bd)$. Hence, whereas the original OFT necessitates $b=d$ to produce a dense orthogonal operator, BOFT can accomplish this using any value of $b$.

\textbf{Identity initialization for BOFT}. Many parameter-efficient finetuning techniques begin from the original pretrained network to limit divergence. For example, LoRA initializes its low-rank matrices with zeros. With BOFT, all butterfly matrices are initialized as identity transformations (meaning that the skew-symmetric matrix in the Cayley parameterization starts as zeros).

\textbf{Multiplicative Dropout for BOFT}. In LoRA~\cite{hulora2022}, Dropout is applied to low-rank update matrices to enhance regularization. Classic Dropout~\cite{srivastava2014dropout} is well suited to LoRA but not immediately applicable for BOFT, which involves multiplicative parameterization. To address this discrepancy, we introduce a specialized multiplicative Dropout for BOFT: since $\bm{R}(m,b)$ is made up of $m$ butterfly matrices, each of which can be rearranged into $2b\times 2b$ block-diagonal form, our dropout method randomly selects a proportion $p_1$ of the butterfly matrices and, within each, a proportion $p_2$ of the diagonal blocks; these are then substituted with identity matrices.

\section{Intriguing Insights and Discussions}

\textbf{Expressivity of BOFT}. While the butterfly structure combined with permutations is capable of exactly reconstructing several established fast linear transforms~\cite{dao2019learning,dao2019kaleidoscope} (\emph{e.g.}, fast Fourier transform, Hadamard transform), it is not yet clear how closely our orthogonal butterfly matrix can approximate a generic orthogonal matrix. To address this, we conduct a simulation where a random dense orthogonal matrix~\cite{anderson1987generation} of dimension $512\times 512$ is approximated. Figure~\ref{fig:para_eff} displays results averaged across 10 different random initializations. On the y-axis, the approximation error is plotted, while the x-axis corresponds to the effective count of trainable parameters. Each curve of a specific color traces BOFT at a fixed block size, with the leftmost marker representing BOFT($1$,$b$) (\emph{i.e.}, the standard block-diagonal OFT with block size $b$). BOFT consistently demonstrates higher parameter efficiency relative to OFT. For instance, BOFT($9$,$2$) surpasses BOFT($1$,$16$) in expressive capability, yet with substantially fewer parameters. BOFT variants with a reduced $b$ and increased $m$ typically attain even better parameter efficiency. An example is BOFT($6$,$4$), which achieves a comparable approximation error to BOFT($2$,$16$) while using significantly fewer parameters. Broadly, the butterfly matrix defines a subset within the orthogonal group possessing more structural constraints (in contrast to block-diagonal forms), underpinning the claim that BOFT exhibits strictly greater expressivity than OFT for the same block size.

\begin{theorem}[Expressivity of BOFT]\label{thm:exp_boft}
    BOFT is strictly more expressive than OFT given identical block sizes. Any orthogonal matrix of dimension $d$ can be decomposed as a product of butterfly matrices:  $\bm{B}_{d-1,1}(d)\bm{B}^\top_{d-1,2}(d)\cdots\bm{B}_{1,1}(d)\bm{B}^\top_{1,2}(d)$, where each $\bm{B}_{i,j}(d)$ for all $i, j$ is a butterfly matrix.
\end{theorem}

Theorem~\ref{thm:exp_boft} inspires a broader framework for BOFT—one where the resulting orthogonal matrix takes the form $\thickmuskip=2mu \medmuskip=2mu \bm{R}^G(m_1,b_1,m_2,b_2,l)=\bm{R}_{l,1}(m_1,b_1)\bm{R}_{l,2}^T(m_2,b_2)\cdots\bm{R}_{1,1}(m_1,b_1)\bm{R}_{1,2}^T(m_2,b_2)$, with $\bm{R}_{i,j}^T(m,b)$ denoting an orthogonal matrix as utilized in BOFT. Setting $m_1=m_2=\log d$, $b_1=b_2=2$, and $l=d-1$ yields $\bm{R}^G(m_1,b_1,m_2,b_2,l)$ capable of representing the complete orthogonal group—a construction recognized as the kaleidoscope hierarchy~\cite{dao2019kaleidoscope}. Nonetheless, we emphasize that increased expressivity does not guarantee improved fine-tuning outcomes; for example, full fine-tuning, while universally expressive, frequently delivers suboptimal performance. Thus, the interplay between expressivity and structural regularity fundamentally determines the transferability of fine-tuned models. BOFT expands the parameter landscape for fine-tuning via the integration of structural priors, equipping practitioners to better negotiate the expressivity-regularity balance.

\textbf{Spectral properties}. Compared to LoRA, orthogonal finetuning typically achieves superior spectral characteristics due to its exact preservation of the spectral norm associated with the pretrained matrix $\bm{W}^0$. This property is readily observed using singular value decomposition: $\bm{W}^0 = \bm{U}\bm{\Sigma}\bm{V}^\top$, with $\bm{U}$ and $\bm{V}$ denoting orthogonal matrices and $\bm{\Sigma}$ containing singular values along its diagonal. In both OFT and BOFT, an orthogonal transformation $\bm{R}$ is applied to the left, producing updated weights $\bm{R}\bm{U}\bm{\Sigma}\bm{V}^\top$, leaving the maximal singular value (and thus the spectral norm of $\bm{W}^0$) unchanged. Maintaining this spectral norm has been demonstrated to substantially enhance stability during training and foster improved generalization~\cite{miyato2018spectral,yoshida2017spectral}. Additional noteworthy mathematical aspects are discussed in Appendix~\ref{app:sn_property}.

\textbf{Orthogonal finetuning as learning bilinear similarity}. The BOFT framework can alternatively be interpreted as optimizing a bilinear similarity function of the form $\bm{w}^0_i\bm{R}\bm{x}$, with $\bm{w}^0_i$ representing the $i$-th column of $\bm{W}^0$. Here, BOFT amounts to finding an orthogonal matrix $\bm{R}$ that encodes the bilinear similarity, subject to stringent constraints—namely, that $\bm{R}$ remains orthogonal—establishing an inherent link to distance metric learning~\cite{xing2002distance} and bilinear forms~\cite{roman2005advanced}.

\textbf{Inductive bias and generalization in BOFT}. The structure of $\bm{R}(m,b)$ in BOFT generally corresponds to a specific subset within the orthogonal group, thereby restricting the hypothesis space and imposing an inductive bias. We posit that the inductive bias imparted by butterfly factorization is advantageous for generalization, owing to its alignment with recurrent structural motifs present in a variety of well-known linear transforms~\cite{dao2019kaleidoscope}, including the discrete Fourier transform, discrete sine/cosine transforms, and the Hadamard transform. In addition, the sparse nature of the matrix factorization employed in BOFT can provide further implicit inductive biases~\cite{gunasekar2017implicit,li2018algorithmic,liu2019neural}.

\textbf{Relation to butterfly-based sparse training}. Several previous studies~\cite{dao2019learning,dao2019kaleidoscope,chen2022pixelated,dao2022monarch} have examined sparse training strategies utilizing butterfly parameterizations. These efforts mainly involve directly reparameterizing neural network weight matrices using butterfly structures and conducting training from the initial state. In contrast, \cite{dao2022monarch} explores finetuning by first projecting pretrained weights onto modified butterfly matrices, then optimizing these projections for specific tasks. The approach proposed by BOFT diverges from these by introducing a unique finetuning process, wherein shared transformation matrices are applied across layers to alter the weights.

\input{rebuttal-results}

\section{Concluding Remarks and Limitations}

This work introduces Orthogonal Butterfly, a broadly applicable parameter-efficient finetuning approach tailored for foundation models, building on butterfly architecture. The central idea for enhanced parameter-efficiency lies in expressing a dense orthogonal matrix as the product of several sparse orthogonal matrices. To facilitate the discovery of suitable matrix decompositions, a graphical framework centered on information transmission is presented. Within this perspective, it becomes evident that butterfly structures serve effectively in enabling sparse orthogonal matrix decompositions. Experimental results substantiate the practical value of BOFT, demonstrating its capacity to efficiently finetune large language models, extensive vision architectures, and text-to-image generation frameworks. The collected empirical evidence underscores the advantage of BOFT as a general-purpose finetuning method.

Nevertheless, there are notable limitations to BOFT. Because the resultant orthogonal transformation arises from multiple sequential orthogonal matrices, the training phase incurs a modestly greater computational overhead compared to OFT. Reducing the training time for BOFT thus remains an unresolved challenge. It is worth noting, however, that post-finetuning, the learned orthogonal transformation can be collapsed into the existing model parameters, so inference speed is unaffected. Additionally, it is still unclear whether the butterfly topology represents the most optimal structure for information transmission. The information transmission framework described here opens the door to insights from areas such as computer networking, which examines optimal topologies for data flow. There is potential for identifying alternative network structures that can further improve the composition of dense orthogonal matrices.